\section{The normal core of a primitive tensor-indecomposable block}
\label{sec:normal-core}

Here tensor indecomposable means that there is no isomorphism
$V\simeq U\otimes W$, with $\dim U,\dim W>1$, under which every
element of the group is a product $a\otimes b$.  A group that
permutes the factors of a tensor decomposition is not excluded by
this hypothesis.  This is the convention preceding
\cite[Theorem~7.1, p.~548]{BI:GS}.

\begin{theorem}[D-BI-19: a forced irreducible quasisimple core]
\label{D-BI-19}
Suppose an inducing block of a rank-nine finite-cover sixfold source
has finite image $J\subset\operatorname{GL}(V)$, is primitive and
tensor indecomposable in the preceding sense, and has degree
$e=\dim V\ge5$.  Write $Z=J\cap\mathbf C^*$, which equals
$Z(J)$ by irreducibility.  Then the following statements hold.
\begin{enumerate}
\item Every nonscalar normal subgroup of $J$ acts irreducibly on
$V$.  Every solvable normal subgroup of $J$ is contained in $Z$.
\item The quotient $J/Z$ has a unique minimal normal subgroup
$S_1\times\cdots\times S_t$, where the $S_i$ are isomorphic
nonabelian simple groups.  Its preimage has the form $ZE$, where
\[
                        E=E_1\cdots E_t                     \tag{BI.69}
\]
is a central product of commuting isomorphic quasisimple groups.
The subgroup $E$ is normal in $J$, acts irreducibly on $V$, and
its factors are permuted transitively by $J$.
\item If the simple groups $S_i$ are of Lie type over
$\mathbf F_q$ and are not alternating, then all inequalities of
D-BI-18 apply with $c(q)$ from (BI.67).  In particular $q\ge47$
is impossible when $e\ge10$, and the exact smaller-field bounds
(BI.64) apply to the total degree $e$.
\item If moreover $e\ge9$, the actual original polarization is
nonsplit, and the representation of $E$ on $V$ is self-dual,
then $q\ge19$ is impossible.  It suffices here that the irreducible
representations on the individual tensor factors are self-dual.
\end{enumerate}
The statement allows larger inducing character fields, outer
automorphisms, and permutations of the quasisimple tensor factors.
\end{theorem}

\begin{proof}
Let $M\triangleleft J$.  Complete reducibility decomposes $V$ into
its $M$-isotypic components.  Normality makes $J$ permute those
components, and irreducibility of $J$ makes the permutation
transitive.  If there were more than one component, they would
give a nontrivial imprimitivity system.  Hence
\[
                  V\simeq W\otimes U,
\]
where $W$ is an irreducible $M$-module and $M$ acts trivially on
the multiplicity factor $U$.  If $M$ is nonscalar, then
$\dim W>1$.  The complex algebra generated by $M$ is
$\operatorname{End}(W)\otimes1_U$, by the matrix-algebra
description of an irreducible summand of the semisimple algebra
$\mathbf C[M]$.

Every element of $J$ normalizes this matrix algebra.  Its induced
automorphism of $\operatorname{End}(W)$ is inner: for example,
the images of a complete set of matrix units furnish a new basis
of the defining module, and change of basis implements the
automorphism.  Thus for each $g\in J$ there is
$a_g\in\operatorname{GL}(W)$ such that
$(a_g^{-1}\otimes1)g$ centralizes
$\operatorname{End}(W)\otimes1$.  The centralizer is
$1\otimes\operatorname{End}(U)$, so
$g=a_g\otimes b_g$ with $b_g\in\operatorname{GL}(U)$.
If $\dim U>1$, this is the forbidden factorwise tensor
decomposition.  Therefore $\dim U=1$, proving that every
nonscalar normal subgroup acts irreducibly.  This argument is
over $\mathbf C$ and uses neither a finite-field specialization
nor maximality of the group.

We next rule out an abelian minimal normal subgroup of $J/Z$.
Let $A/Z$ be such a subgroup and let $A$ be its full preimage in
$J$.  It is nontrivial in $J/Z$, so $A$ is nonscalar and acts
irreducibly by the preceding paragraph.  Since $A/Z$ is abelian,
all commutators in $A$ lie in the scalar group $Z$.  A nonscalar
$x\in A$ is not central in $A$, by Schur's lemma, so some
$h\in A$ has $hxh^{-1}=zx$ with a nonidentity scalar $z$.
Taking traces gives $\operatorname{tr}(x)=z\operatorname{tr}(x)$,
and hence every nonscalar $x\in A$ has trace zero.  D-BI-18 with
$N=A$ and $c=0$ then forces $e\in\{3,4\}$, contradicting
$e\ge5$.  This excludes the entire abelian minimal-normal case,
without needing an extraspecial-group classification.

The quotient $J/Z$ is nontrivial, since $e>1$ and $J$ is
irreducible.  Choose any minimal nontrivial normal subgroup
$\overline N$ of $J/Z$.  A minimal normal subgroup of a finite
group is characteristically simple and is therefore a direct
product of isomorphic simple groups.  For completeness, choose a
minimal nontrivial normal subgroup $T$ of $\overline N$ and take
all its distinct images under automorphisms of $\overline N$.
They generate $\overline N$ by characteristic simplicity.
Distinct such subgroups commute, since their intersection is
trivial by minimality.  Every normal subgroup of $T$ is then
normal in the generated group, so $T$ itself is simple.  If $T$
is abelian it has prime order and the full product is elementary
abelian; otherwise its center is trivial and the product is a
direct product of nonabelian simple groups.  The abelian case has
just been excluded.  We therefore have
\[
                \overline N=S_1\times\cdots\times S_t
\]
with the $S_i$ isomorphic and nonabelian simple.  The action of
$J/Z$ on these factors is transitive: the product over any orbit
is a normal subgroup inside $\overline N$, so minimality leaves
only one orbit.

Let $N$ be the preimage of $\overline N$, and let $A_i$ be the
preimage in $N$ of $S_i$.  As $A_i/Z$ is perfect,
$A_i=A_i'Z$.  Consequently
\[
           A_i'=[A_i'Z,A_i'Z]=[A_i',A_i'],
\]
so $E_i:=A_i'$ is perfect.  Its map onto $S_i$ has central
kernel $E_i\cap Z$, and its center is exactly that kernel,
because $S_i$ has trivial center.  Thus $E_i$ is quasisimple.
For $i\ne j$, the commutators $[A_i,A_j]$ are scalar.  With
$y\in A_j$ fixed, $x\mapsto[x,y]$ is then a homomorphism from
$A_i$ to the abelian group $Z$, so it vanishes on $E_i$.
It follows that $[E_i,A_j]=1$, and hence that $E_i$ and $E_j$
commute.  Since each $A_i=E_iZ$ and the $A_i/Z$ generate
$N/Z$, we obtain
\[
                        N=Z E_1\cdots E_t,
                        \qquad N'=E_1\cdots E_t=E.
\]
The intersections of distinct factors are central, proving
(BI.69).  Conjugation by $J$ permutes the $A_i$ and therefore
the $E_i$ transitively.  In particular the factors are isomorphic.
The subgroup $E=N'$ is normal in $J$, and it acts irreducibly:
$N$ is nonscalar normal and irreducible, and adding the scalar
group $Z$ does not change invariant subspaces.

To see uniqueness of $\overline N$, suppose
$\overline M$ were a different minimal normal subgroup of
$J/Z$.  Distinct minimal normal subgroups commute.  If $M$ is
its preimage, this gives $[M,E]\subset Z$.  For any fixed
$m\in M$, the map $x\mapsto[m,x]$ from the perfect group $E$
to $Z$ is a homomorphism, hence trivial.  Thus $M$ centralizes
the irreducible group $E$, so Schur's lemma implies $M\subset Z$,
a contradiction.  Similarly, a nontrivial solvable normal subgroup
of $J/Z$ would contain a minimal normal subgroup of $J/Z$ that
is elementary abelian.  That case was excluded above.  As $Z$
is solvable and central, every solvable normal subgroup of $J$
is contained in $Z$.

The representation of the commuting central product $E$ is a
tensor product of irreducible representations of the $E_i$.
Each factor representation is faithful: an element acting as the
identity there would act as the identity on $V$, and $J$ is an
actual matrix group.  If the $S_i$ are of Lie type over
$\mathbf F_q$, the published bound (BI.67) applies to each
quasisimple factor.  The central-product assertion of D-BI-18
therefore bounds the normalized character of the irreducible
normal group $E$ on $V$ by the same $c(q)$.  Its degree in
(BI.62) is the total degree $e$, not a single tensor-factor
degree.  The quoted consequences now follow from D-BI-18.

If the representation of $E$ is self-dual, its invariant
nondegenerate bilinear form is unique up to scalar, so its
normalizer $J$ is essentially self-dual as proved in D-BI-18a.
If all the tensor factors are self-dual, tensoring their
nondegenerate forms supplies such a form on $V$.  The final
assertion follows from D-BI-18a, which allows arbitrary inducing
character fields and invokes the separately proved D-BI-09.
\end{proof}

\paragraph{Relation to the published structure reduction.}
The homogeneous-restriction, tensor-normalizer, and central-product
steps have the same form as the proof of
\cite[Theorem~5.1, pp.~148--151]{BI:DHP}.  That theorem itself
assumes a maximal coprime group over a finite field.  We have not
applied that theorem outside its hypotheses: the complete
characteristic-zero argument needed here is written above.
The result imposes no field-size bound on alternating or sporadic
simple factors and leaves the indicated small Lie-type fields
open.  A primitive group preserving a nontrivial factorwise
tensor decomposition is also outside this theorem.
